--- Page 1 ---
 
1 
Supplementary Information 
 
Combining Valence-to-Core X-ray Emission and Cu K-edge X-ray Absorption 
Spectroscopies to Experimentally Assess Oxidation State in Organometallic 
Cu(I)/(II)/(III) Complexes 
 
Blaise L. Geoghegan,a,b Yang Liu,c Sergey Peredkov,a Sebastian Dechert,c Franc Meyer,c Serena DeBeer,a 
George E. Cutsail IIIa,b* 
a Max Planck Institute for Chemical Energy Conversion, Stiftstrasse 34–36, 45470 Mülheim an der Ruhr, 
Germany 
b Institute of Inorganic Chemistry, University of Duisburg-Essen, Universitätsstrasse 5–7, 45117 Essen, Germany 
c Institute of Inorganic Chemistry, University of Göttingen Tammannstrasse 4, 37077 Göttingen, Germany 
 
email: george.cutsail@cec.mpg.de 
 
 


--- Page 2 ---
 
2 
Contents 
Contents .......................................................................................................................................................... 2 
Experimental ................................................................................................................................................... 3 
Sample Preparation and Synthesis ............................................................................................................. 3 
X-Ray Structure Determination .................................................................................................................... 4 
X-ray Absorption Data Collection ................................................................................................................ 4 
X-ray Emission Data Collection ................................................................................................................... 4 
Density Functional Theory ........................................................................................................................... 5 
NMR Spectroscopy ...................................................................................................................................... 6 
UV-Vis and infrared spectroscopy ............................................................................................................. 13 
Mass Spectrometry .................................................................................................................................... 14 
SQUID magnetometry ............................................................................................................................... 14 
X-ray Crystallography ................................................................................................................................ 16 
X-ray Spectroscopy ....................................................................................................................................... 18 
X-ray Absorption ........................................................................................................................................ 18 
X-ray Emission ........................................................................................................................................... 20 
Cu Kβ Mainlines ..................................................................................................................................... 20 
Cu Kβ VtC XES ...................................................................................................................................... 21 
UV-Vis ........................................................................................................................................................ 30 
Example Orca Input Files .......................................................................................................................... 32 
Geometry Optimization .......................................................................................................................... 32 
X-ray Absorption Spectroscopy ............................................................................................................. 32 
X-ray Emission Spectroscopy ................................................................................................................ 32 
Löwdin and Mulliken Population analysis of QROs ............................................................................... 33 
Broken Symmetry ................................................................................................................................... 33 
UV-vis spectroscopy .............................................................................................................................. 33 
XYZ Coordinates for Orca Calculations ..................................................................................................... 34 
References .................................................................................................................................................... 40 
 
 
 


--- Page 3 ---
 
3 
Experimental 
Sample Preparation and Synthesis 
General Considerations. CuCl, CuF2, CuCl2 were purchased from Strem Chemicals and used as received. 
CF3SiMe3, and KF were purchased from Sigma-Aldrich (Germany) and used as received. [CuI(NHC2)](PF6), 
[CuII(NHC2)](PF6)2∙CH3CN and [CuIII(NHC2)](PF6)2(SbF6)∙CH3CN were prepared as previously described.1 
[NBu4][Cu(CF3)4] was synthesized via the method described previously by A. M. Romine et al.2 The 
tetraimidazolium salt calix[4](4,5-dimethylimidazolium) hexafluorophosphate ([H4LMe](PF6)4) was synthesized 
according to reported literature procedures.3 Me3CN-d3 was dried with 3 Å molecular sieves and stored in the 
glovebox. All other chemicals and solvents were purchased from commercial suppliers, and used without further 
purification. All manipulations of air- and moisture-sensitive materials were performed under an atmosphere of 
dry dinitrogen with the rigorous exclusion of air and moisture using standard Schlenk techniques, or in a N2-filled 
glovebox. 1H, 31P and 19F NMR spectra were recorded with Bruker 300 NMR spectrometer at 300 MHz, 122 MHz 
and 282 MHz, respectively. 13C NMR spectra were recorded with Bruker 500 NMR spectrometer at 125 MHz. 
Variable-temperature NMR were recorded with Bruker 400 NMR spectrometer at 400 MHz. All NMR chemical 
shifts were reported in units of ppm with references to the residual protons of the deuterated solvents for proton 
chemical shifts (MeCN-d3 δH = 1.94 ppm), the 13C of deuterated solvents for carbon chemical shifts (MeCN-d3 
δC = 1.32, 118.36 ppm), the 19F of CFCl3 (external standard) for fluorine chemical shifts, and the 31P of 85% 
phosphorous acid (external standard) for phosphine chemical shifts. ESI mass spectra were recorded on a Bruker 
HCT ultra spectrometer. Elemental analyses were performed by the analytical laboratory of the Institute of 
Inorganic Chemistry at the University of Göttingen using an Elementar Vario EL III instrument. UV-Vis spectra 
were recorded with an Agilent Cary 60. IR spectra were recorded inside a glovebox on a Cary 630 FTIR 
spectrometer equipped with  a Diamond ATR module and analyzed by FTIR MicroLab software. 
Preparation of [Cu(NHC4)](PF6)3. This complex was synthesized via a procedure described in literature for 
the parent non-methylated complex,4 with slight modifications. [H4LMe](PF6)4 (351 mg, 0.35 mmol) and copper(II) 
acetate monohydrate (138 mg, 0.69 mmol) were dissolved in 10 mL of reagent grade dimethyl sulfoxide in air. 
The blue reaction mixture was stirred at room temperature for 30 min and then stirred at 40 oC for 5 h. During 
this time the color of the solution changed from blue to clear green. The solution was cooled to room temperature 
and the crude product was precipitated with a large amount of dichloromethane (70 mL), forming a fine white 
precipitate. After filtration, the resulting white solid was consecutively washed with dichloromethane (3 mL × 3) 
and Et2O (5 mL × 3) and dried under vacuum to give [Cu(NHC4)](PF6)3 as a white powder (215 mg, 67% yield). 
Yellowish crystals suitable for X-ray diffraction studies were grown by diffusing Et2O into a saturated MeCN 
solution of the complex at room temperature. 1H NMR (300 MHz, MeCN-d3, 298 K): δ (ppm) 6.12 (s, 8H, CH2-
linkers), 2.49 (s, 24H, -CH3). 13C NMR (125 MHz, MeCN-d3, 298 K): δ (ppm) 149.5 (NCN), 129.7 (C-imidazolyl 
backbone), 59.9 (CH2-linkers), 9.32 (-CH3). 19F NMR (282 MHz, MeCN-d3, 298 K): δ (ppm) -72.9 (d, JP-F = 707.2 
Hz). 31P NMR (122 MHz, MeCN-d3, 298 K): δ (ppm) -144.7 (hept, JF-P = 707.2 Hz). Anal. Calcd. for 
CuC24H32N8P3F18: C 30.96, H 3.46, N 12.04; Found: C 30.72, H 3.38, N 11.96. ESI–MS (MeCN) m/z (%): 785.1 
(100) [M-PF6]+. Absorption spectrum (in MeCN): λmax, nm (ε, M-1 cm-1) 230 (27200), 260 (34200), 347 (1740). 
ATR-IR (powder, cm−1): ν = 1641 (w), 1496 (m), 1440 (w), 1405 (w), 1381 (m), 1266 (w), 1199 (w), 1044 (w), 944 
(w), 921 (w), 822 (s), 739 (m), 711 (w), 649 (w), 555 (s), 531 (m), 519 (w). 


--- Page 4 ---
 
4 
X-ray spectroscopy sample preparation. For XAS experiments, samples were diluted with boron nitride to 
achieve a 2% (w/w) concentration of copper, then packed into 1 mm thick aluminum sample cells and sealed with 
38 μm Kapton tape. For XES experiments, all samples were measured in the solid state at room temperature. 
The pure solids were ground to a fine powder and packed into 0.5 mm thick aluminum sample holders and sealed 
with 38 μm Kapton tape. 
X-Ray Structure Determination 
Crystal data and details of the data collections for [Cu(NHC4)](PF6)3 are given in Table S2, the molecular 
structure of the cation [CuIII(NHC4)]3+ is shown in Figure S16. X-ray data were collected on a BRUKER D8-QUEST 
diffractometer (monochromated Mo-Kα radiation, λ = 0.71073 Å) by use of w and f scans at low temperature. 
The structure was solved with SHELXT and refined on F2 using all reflections with SHELXL-2018.5 Non-hydrogen 
atoms were refined anisotropically. Hydrogen atoms were placed in calculated positions and assigned to an 
isotropic displacement parameter of 1.5/1.2 Ueq(C). Four F-atoms of one PF6– counterion were found to be 
disordered (occupancy factors 0.48(3)/0.52(3)). SADI (dF∙∙∙F) restrain were applied to model the disordered parts. 
Absorption correction was performed by the multi-scan method with SADABS.6 CCDC 2104793 contains the 
supplementary crystallographic data for this paper. These data can be obtained free of charge from the 
Cambridge Crystallographic Data Centre via www.ccdc.cam.ac.uk/structures. 
X-ray Absorption Data Collection 
Data collection procedures for samples of [Cun+(NHC2)]n+ complexes were previously described1 albeit, only 
the Kβ high-energy resolution fluorescence detected (HERFD) XAS was reported. Simultaneous transmission 
data was also collected at the Stanford Synchrotron Radiation Lightsource (SSRL), BL6-2, and is now reported 
here. Comparisons of the Kβ-HERFD and transmission XAS are shown in Figure S17. [CuIII(NHC4)]3+ and 
[CuIII(CF3)4]− were measured in transmission mode at room temperature and ambient conditions at the SAMBA 
beamline of the SOLEIL synchrotron under ring conditions of 3 GeV and 500 mA. Incident X-ray radiation was 
monochromatized by using a Si(220) 
 sagittal focusing monochromator and calibrated by using a Cu foil 
calibration standard and assigning the energy of the first inflection point to 8980.3 eV. Five scans were measured 
and averaged using in-house developed MATLAB R2019b scripts and all further data processing was conducted 
using the Athena program for x-ray spectroscopic data processing.7 The background was accounted for by fitting 
a polynomial to the pre-edge region and subtracting this polynomial from the entire spectrum. All spectra were 
normalized by fitting a flattened polynomial to the post-edge region of the spectra and normalizing the edge jump 
to 1.0. No spectral changes due to photodamage were observed after ten 65 second scans on a single sample 
spot (beam size = 2000 x 300 μm {width x height}). 
X-ray Emission Data Collection 
XES data collection was done at the PINK tender X-ray beamline at BESSY II. The PINK beamline is currently 
being operated in commissioning mode. A considerable gain in fluorescence signal could be obtained by using 
a multilayer monochromator (∼100 eV band pass), and the energy of the incident beam was 9.5 keV. The beam 
size on the sample was 30 μm × 500 μm (V×H). All spectra were collected using an in-house-designed energy-
dispersive von Hamos spectrometer with a vacuum sample chamber environment (5 mbar working pressure). 


--- Page 5 ---
 
5 
The analyser was set up in a vertical dispersion direction, taking advantage of the small vertical beam size to 
improve the energy resolution. A Si(444) crystal with a bending radius of R=250mm dispersed incoming 
fluorescence radiation onto an Eiger detector with a 75 μm × 75 μm pixel size. The 80mm x 40mm detector chip 
accepted fluorescent radiation reflected from the crystal analyser under 68º-60.7º brag angles that corresponded 
to an energy window of 8530-9065 eV. Calibration of the energy scale was achieved by measuring reference X-
ray emission lines of anhydrous CuCl in the same instrumental configuration and setting the energy of the Kβ1,3 
mainline feature to 8904.7 eV and VtC feature to 8977.4 eV (Figure S1). We have previously collected the Kβ 
emission of anhydrous CuCl at both ID26, ESRF, Grenoble, France, Figure 20. The spectrometer at ID26 was 
calibrated by scanning of elastic scattering lines through the available emission range of the spectrometer. The 
monochromator for each was calibrated by setting the first inflection point of a Cu foil to 8980.3 eV. 
Cu Kβ XES damage scans and assessments were performed for all samples. In all cases, the data were 
collected in continuous motion. Samples were scanned at a rate of 100 µm/s, resulting in a total exposure time 
of 0.50 s/pass and the total exposure was between 6 and 10 s per sample. None of the samples exhibited damage 
or changes in the Kβ emission spectra within the measurement scan periods. XES data were processed and fit 
using in-house developed MATLAB scripts and the integrated intensity of the Kβ main line was set to 1.0. 
 
Figure S1 Cu Kβ mainline (left) and VtC region (right) for solid CuCl energy calibrants as measured at the ID26 
beamline at ESRF, Grenoble, and the PINK beamline at BESSY II, Berlin. The spectra were normalized by setting 
the area of the mainline region (8860-8930 eV) to 1.0 a.u. 
 
Density Functional Theory 
All DFT and TDDFT were performed with the ORCA electronic structure package 4.2.8,9 Geometry 
optimization calculations were carried out at the B3LYP10,11 level of theory, using def2- variants of Ahlrichs’ all-
electron Gaussian triple-ζ valence polarized recontracted basis set (def2-TZVP)12 on all atoms and the AutoAux 
basis option for ORCA.13 The calculations employed the resolution of identity (RI-J) algorithm for the computation 
of the Coulomb terms and the ‘chain of spheres exchange’ (COSX) algorithm for the calculation of the exchange 
terms14 and a tight self-consistent field (SCF) convergence threshold was chosen via the “TightSCF” keyword. 
An increased grid was used during the SCF iterations (Grid4), and for the final energy evaluation after SCF 


--- Page 6 ---
 
6 
convergence Grid7 was used. The conductor-like polarizable continuum model (CPCM) was used for charge 
compensation in all calculations of complexes carrying a net positive/negative charge.15 Vibrational frequencies 
were calculated for all optimized structures and the absence of imaginary modes alongside six 0 cm-1 modes 
confirmed that the true minima were obtained in all cases.  
Cu Kβ VtC XES were calculated using the features described above, with the exception that a scalar 
relativistic basis set (ZORA-def2-TZVP)12 was employed. In the XES block, “CoreOrb” dictates the orbitals that 
will be the electron acceptors. For transition metal XES spectra, the metal 1s orbital is usually orbital 0, and is 
thus selected. OrbOp defines the operators of the electrons (α = 0, β = 1) that will be calculated, both of which 
were selected. CoreOrbSOC defines which core orbitals (α = 0, β = 1) are treated with spin-orbit coupling (SOC), 
both of which were selected to accommodate the OrbOp selection. A 2.0 eV full-width half-max Gaussian 
broadening and -5.1 eV energy shift was applied when plotting all calculated XES spectra.  
Time-dependent DFT (TD-DFT) calculations used the B3LYP functional, ZORA-def2-TZVP basis set, 
AutoAux and the zeroth-order regular approximation (ZORA) to account for scalar relativistic effects.16,17 TDDFT 
calculated XAS were shifted by -3.1 eV and a 2.5 eV full-width half-max Gaussian broadening was applied when 
plotting all calculated XES spectra. UV-vis spectra were broadened by 2,500 cm-1. 
NMR Spectroscopy 
 
Figure S2 1H NMR spectrum of [Cu(NHC4)](PF6)3 in MeCN-d3 (300 MHz, 298 K). 


--- Page 7 ---
 
7 
 
Figure S3 13C NMR spectrum of [Cu(NHC4)](PF6)3 in MeCN-d3 (125 MHz, 298 K). 
 
Figure S4 19F NMR spectrum of [Cu(NHC4)](PF6)3 in MeCN-d3 (282 MHz, 298 K). 


--- Page 8 ---
 
8 
 
Figure S5 31P NMR spectrum of [Cu(NHC4)](PF6)3 in MeCN-d3 (122 MHz, 298 K). 
 
Figure S6 1H-1H NOESY spectrum of [Cu(NHC4)](PF6)3 in MeCN-d3 (300, 300 MHz, 298 K). 


--- Page 9 ---
 
9 
 
Figure S7 1H-13C HSQC spectrum of [Cu(NHC4)](PF6)3 in MeCN-d3 (300, 75 MHz, 298 K). 
 
Figure S8 1H-13C HMBC spectrum of [Cu(NHC4)](PF6)3 in MeCN-d3 (300, 75 MHz, 298 K). 
 
 


--- Page 10 ---
 
10 
VT-NMR and Dynamic Behavior 
At room temperature, the 1H NMR spectrum of [Cu(NHC4)](PF6)3 in MeCN-d3 shows two sharp singlets at 
6.12 ppm (CH2 linkers) and 2.49 ppm (CH3) (Figure S9). Upon decreasing the temperature, signal broadening 
followed by splitting of the peak at 6.12 ppm (CH2 linkers) into a pair of apparent doublets (AB system, Jab = 13.6 
Hz) is observed (Figures S8). This suggests that the dynamic behavior (conformational inversion of the 
macrocyclic ligand) becomes slow on the NMR time scale at low temperatures (coalescence temperature Tc = 
270 K, ΔG‡ = 12.7 kcal∙mol−1). Further analyses of the VT-NMR data yielded the activation energy (Ea = 18.3 
kcal∙mol–1) as well as activation parameters (ΔH‡ = 17.8 kcal∙mol–1 and ΔS‡ = 19.0 cal∙mol–1∙K–1) from Arrhenius 
and Eyring plots, respectively (Table S1, Figures S9 and S10). 
 
Figure S9 Variable-temperature 1H NMR spectra of [Cu(NHC4)](PF6)3 in MeCN-d3 (400 MHz, temperature range 
from 233 K to 313 K). Only regions 5.0–7.0 ppm (the intensity was scaled up by a factor of 50) and 2.0–3.0 ppm 
are shown for clarity.  
To determine these activation parameters for the ring inversion of [Cu(NHC4)](PF6)3, the rate constants were 
calculated according to the formulas in literature.18 The line width at half height was determined by using the line 
fitting function of MestReNova program. Ring inversion was assumed to be frozen at 233 K in this case. The 
Gibbs free energy of activation Δ‡G at coalescence temperature Tc can be calculated from equation (1). The 
calculation of thermodynamic parameters (activation energy, enthalpy and entropy) was done according to the 
Arrhenius equation (2) and Eyring equation (3). 
𝛥𝐺‡ = 𝑅𝑇" (𝐼𝑛
#!
$! + 𝐼𝑛
$"
% ,          (1) 


--- Page 11 ---
 
11 
𝐼𝑛 𝑘 = −
&#
'
(
# + 𝐼𝑛(𝐴)                  (2) 
𝐼𝑛
$
#
 
= − 
*+‡
'
(
# +
*,‡
' + 𝐼𝑛
$%
%          (3) 
Table S1 Calculated values for k at different temperatures T. 
 
T (K) 
k (s-1) 
Ink (s-1) 
In(k/T) (s-1 K-1) 
238 
3.2 
1.155 
-4.318 
243 
7.8 
2.057 
-3.436 
248 
16.1 
2.780 
-2.734 
253 
32.2 
3.472 
-2.061 
258 
78.2 
4.360 
-1.193 
263 
162.8 
5.092 
-0.480 
268 
272.4 
5.607 
0.016 
270 
323.8 
5.780 
0.182 
273 
424.4 
6.051 
0.441 
278 
806.2 
6.692 
1.065 
283 
1733.0 
7.458 
1.812 
298 
8948.4 
9.099 
3.402 
303 
12937.1 
9.468 
3.754 
 
Figure S10 Arrhenius plot for the determination of Ea for the ring inversion in [Cu(NHC4)](PF6)3. The following 
value was determined: Ea = 18.3 ± 0.2 kcal∙mol-1. 


--- Page 12 ---
 
12 
 
Figure S11 Eyring-plot for the determination of ∆H‡ and ∆S‡ for the ring inversion in [Cu(NHC4)](PF6)3. The 
following values were determined: ∆H‡ = 17.8 ± 0.2 kcal∙mol-1, ∆S‡ = 19.0 ± 0.8 cal∙mol-1 K-1. 
 
 
 


--- Page 13 ---
 
13 
UV-Vis and infrared spectroscopy 
 
Figure S12 UV-Vis spectra of [Cu(NHC4)](PF6)3  (in MeCN) recorded at room temperature. 
 
Figure S13 ATR-IR spectrum of solid [Cu(NHC4)](PF6)3. 
 
 


--- Page 14 ---
 
14 
Mass Spectrometry 
 
Figure S14 ESI–MS of complex [Cu(NHC4)](PF6)3  in MeCN. The inset shows the experimental (bottom) and 
simulated (top) isotope pattern for the peak around m/z = 785.1 (100) [M-PF6]+. 
SQUID magnetometry 
Temperature-dependent magnetic susceptibility measurements for [CuIII(NHC4)](PF6)3 
were carried out with a Quantum-Design MPMS-3 SQUID magnetometer equipped with a 7.0 
T magnet in the range from 2-350 K on a polycrystalline powdered sample under an applied 
magnetic field of 0.5 T. The crystalline solid sample was contained in a polycarbonate capsule 
and fixed in a nonmagnetic sample holder. The raw data file for the measured magnetic 
moment was corrected for the diamagnetic contribution of the capsule according to Mdia = χg 
× m × H, with experimentally obtained gram susceptibilities of the capsule. The molar 
susceptibility data of the compound were corrected for the diamagnetic contribution. 
Experimental data for [Cu(NHC4)](PF6)3  were modelled with the julX program.19 


--- Page 15 ---
 
15 
 
Figure S15 χMT vs. T measurement for a solid sample of [CuIII(NHC4)](PF6)3  in the temperature range of 2-350 
K at 0.5 T. The data shows that [CuIII(NHC4)](PF6)3  is diamagnetic over the entire temperature range studied. 
 
 


--- Page 16 ---
 
16 
X-ray Crystallography 
Table S2 Crystal data and structure refinement of [Cu(NHC4)](PF6)3·2CH3CN 
compound 
1 (yl35) 
empirical formula 
C28H38CuF18N10P3 
moiety formula 
C24H32CuN83+, 3(F6P–), 2(C2H3N) 
formula weight 
1013.13 
T [K] 
100(2) 
crystal size [mm³] 
0.380 x 0.126 x 0.071 
crystal system 
orthorhombic 
space group 
Pca21 (No. 29) 
a [Å] 
18.4087(9) 
b [Å] 
12.9782(6) 
c [Å] 
16.4884(8) 
V [ų] 
3939.3(3) 
Z 
4 
r [g·cm–³] 
1.708 
F(000) 
2048 
µ [mm–1] 
0.800 
Tmin / Tmax 
0.85 / 0.94 
q–range [°] 
2.213 – 27.895 
hkl–range 
–23 to 24, ±17, ±21 
measured refl. 
37226 
unique refl. [Rint] 
9387 [0.0376] 
observed refl. (I > 2s(I)) 
8167 
data / restr. / param. 
9387 / 29 / 587 
goodness–of–fit (F²) 
1.023 
R1, wR2 (I > 2s(I)) 
0.0311 / 0.0750 
R1, wR2 (all data) 
0.0390 / 0.0804 
res. el. dens. [e·Å–³] 
–0.381 / 0.574 
 


--- Page 17 ---
 
17 
 
Figure S16 A) Crystal structure of the complex cation [CuIII(NHC4)]3+ as viewed down the molecular z-axis. B)  
View of the complex cation down one of its dihedral mirror planes (σd1). C) View of the complex cation down the 
second of its dihedral mirror planes (σd2). Acetonitrile solvates and PF6− anions are omitted for clarity. Ellipsoids 
are shown at the 50% probability level. Legend: copper, orange; carbon, grey; nitrogen, blue. 
 
 
Table S3 Selected bond length and angles for the complex cation [CuIII(NHC4)]3+. 
Bonds 
[CuIII(NHC4)]3+ 
Bond lengths [Å] 
 
Cu1-C1 
1.877(3) 
Cu1-C7 
1.881(3) 
Cu1-C13 
1.889(3) 
Cu1-C19 
1.881(3) 
Av. Cu-C 
1.882(3) 
 
 
Bond angles [°] 
 
C1-Cu1-C13 
178.90(14) 
C7-Cu1-C19 
178.11(14) 
C1-Cu1-C7 
89.60(14) 
C1-Cu1-C19 
89.57(13) 
C7-Cu1-C13 
90.56(13) 
C13-Cu1-C19 
90.25(14) 
Av. trans C-Cu-C 
178.51(14) 
Av. cis C-Cu-C 
90.00(14) 
 


--- Page 18 ---
 
18 
X-ray Spectroscopy 
X-ray Absorption 
 
Figure S17 Comparison of the Kβ high energy resolution fluorescence detected (HERFD)1 and Transmission 
mode Cu K-edge XAS for the [Cun+(NHC2)]n+ series. The inset shows the magnification of the pre-edge region 
for the Cu(II) and Cu(III) species. 
 


--- Page 19 ---
 
19 
 
Figure S18 B3LYP/def2-TZVP calculated LUMOs for Cu(III) complexes of the NHC2 (left) and NHC4 (right) 
macrocyclic ligands. Isosurfaces are plotted at 0.05 a.u. 
 
Table S4 Selected orbital energies and corresponding Cu d/p orbital character for [CuIII(NHC4)]3+ and 
[CuIII(CF3)4]3+. Orbital compositions were determined from reduced Löwdin orbital population analysis. 
 
[CuIII(NHC4)]3+ 
[CuIII(CF3)4]− 
Orbital 
Energy [eV] 
Cu nd orbital [%]  
Cu np orbital [%] 
Energy [eV] 
Cu nd orbital [%]  
Cu np orbital [%] 
LUMO +2 
-0.86 
0.7 xz, 7.6 yz 
0 
0.97 
11.4 z2 
0 
LUMO +1 
-2.25 
0 
15.1 pz 
0.23 
3.2 x2-y2 
32.2 pz, 0.9 py 
LUMO 
-3.29 
43.4 x2-y2 
0 
-2.64 
35.5 x2-y2 
3.5 pz 
HOMO 
-6.97 
0 
0 
-7.51 
1.2 xy, 0.8 yz 
5.6 px, 6.6 py 
HOMO-1 
-7.02 
2.9 xy, 0.9 yz 
0 
-7.53 
1.0 xy, 1.1 yz 
6.7 px, 5.3 py 
HOMO-2 
-7.02 
0 
2.7 pz 
-8.83 
34.8 z2 
0 
HOMO-3 
-7.44 
0 
0 
-10.03 
80.4 yz 
0 
HOMO-8 
-9.67 
0.1 xz, 0.1 yz 
5.7 px, 6.7 py 
-10.75 
0 
0 
HOMO-9 
-9.68 
0.1 xz, 0.1 yz 
6.7 px, 5.7 py 
-10.87 
0 
0 
 
 
 
 


--- Page 20 ---
 
20 
 
Figure S19 Experimental (top) and TDDFT calculated (bottom) XAS for the three Cu(III) complexes investigated 
along with transition difference densities for selected pre-edge and rising edge transitions. transition difference 
densities are plotted with an isosurface value of 0.005 a.u. Calculated spectra are shifted -3.1 eV. 
X-ray Emission 
Cu Kβ Mainlines 
Table S5 Energies of the Kβʹ and Kβ1,3 features, and their separation, for all complexes considered in the present 
study. 
 
Kβ’ (eV) 
Kβ1,3 (eV) 
ΔE (eV) 
CuCl 
8895.2 
8904.6 
9.4 
CuCl2 
8894.7 
8904.7 
9.9 
CuF2 
8896.5 
8904.7 
8.2 
[CuI(NHC2)]+ 
8895.3 
8904.8 
9.5 
[CuII(NHC2)]2+ 
8895.1 
8904.9 
9.8 
[CuIII(NHC2)]3+ 
8895.7 
8904.9 
9.2 
[CuIII(NHC4)]3+ 
8895.1 
8904.9 
9.8 
[CuIII(CF3)4]− 
8895.7 
8904.9 
9.2 
 


--- Page 21 ---
 
21 
Cu Kβ VtC XES 
 
Figure S20 a) Background subtracted VtC spectra for the [Cun+(NHC2)]n+ series and [CuIII(NHC4)]3+ and b) 
Background subtracted VtC spectra for the Cu(III) complexes [CuIII(NHC2)]3+, [CuIII(NHC4)]3+ and [CuIII(CF3)4)]− 
with individual fitted Gaussian bands (colored) and their sum (red). Residual intensity (experiment – fit) are shown 
in black below the individual VtC spectra.  
 
Table S6 Peak analysis for the single dominant VtC emission transition in all complexes. 
Compound 
Energy [eV] 
FWHM [eV] 
Integrated Intensitya 
[CuI(NHC2)]+ 
8975.9 
1.59 
5.0 
[CuII(NHC2)]2+ 
8977.9 
1.23 
3.1 
[CuIII(NHC2)]3+ 
8979.0 
1.51 
4.4 
[CuIII(NHC4)]3+ 
8979.0 
1.21 
5.9 
[CuIII(CF3)4)]− 
8979.9 
1.30 
5.0 
a Integrated intensities were multiplied by 1000. 
  


--- Page 22 ---
 
22 
 
Figure S21 Overlay of the high energy Cu Kβ VtC XES region (dashed) and the edge region of the Cu K-edge 
XAS (solid) for [CuII(NHC2)]2+, [CuIII(NHC2)]3+ and [CuIII(NHC4)]3+. 
 
 
Figure S22 Crystal structures of the cationic units of the [Cun+(NHC2)]n+ oxidation series. Legend: carbon, gray; 
copper, orange; nitrogen, blue. Non-coordinated solvates, anion and hydrogen atoms are omitted for clarity. 
Thermal ellipsoids are shown at the 50% probability level. 
 
Table S7 Single point energies for the singlet and triplet state in Cu(III) complexes relative to their minimum 
energy solution. 
 
[CuIII(NHC2)]3+ 
[CuIII(NHC4)]3+ 
[Cu(CF3)4]− 
Singlet energy [eV] 
 
 
 
UKS  
0 
0 
0 
RKS 
0 
0 
0 
BS (MS = 0)a 
N/A 
N/A 
N/A 
Triplet energy [eV] 
 
 
 
UKSb 
+1.5 
+3.0 
+1.4 
a Broken symmetry DFT single point calculation was conducted but no solution was found. b Triplet state energy 
calculated via a single point calculation using the xyz file obtained from the geometry optimization of the Cu(III) 
complex in the triplet state. 


--- Page 23 ---
 
23 
 
 
Figure S23 Energies of the calculated vs. experimental emission maxima in the VtC region of the Cu Kβ XES 
for all complexes. The red line is a linear fit of the data with the slope fixed to a value of 1.000 to extract a 
constant energy shift from the y intercept. 
 


--- Page 24 ---
 
24 
 
Figure S24 Calculated VtC spectra for [CuI(NHC2)]+, [CuII(NHC2)]2+ and [CuIII(NHC2)]3+ with donor MOs 
corresponding to the most intense one-electron transitions. Isosurfaces are plotted at a value of 0.05 a.u. Spectra 
are shifted by -5.1 eV and a linewidth broadening of 2.0 eV was applied. 
 


--- Page 25 ---
 
25 
Table S8 Calculated orbital contributions to the donor MOs responsible for the main Kβ2,5 feature in the VtC 
spectra. Values are obtained from both α and β orbital sets and the VtC region was determined by summing the 
intensities of transitions encompassed by a gaussian line function fit to the selected energy region of the VtC 
spectrum. 
 
% Cu np 
% Cu nd 
% ligand np 
% ligand ns 
[CuI(NHC2)]+ 
7.3 
8.6 
64.8 
11.1 
[CuII(NHC2)]2+ 
4.4 (2.7)
a 
14.4 (15.6)
a 
63.5 (64.4)
a 
9.7 (8.3)
a 
[CuIII(NHC2)]3+ 
3.5 
6.6 
70.0 
12.2 
[CuIII(NHC4)]3+ 
9.2 
1.8 
65.8 
15.5 
[CuIII(CF3)4]– 
10.1 
5.3 
71.2 
8.8 
a In open shell complexes, values outside of brackets are given for the α spin manifold, values given inside 
brackets are for the β spin manifold. MO range for the accumulative Löwdin population analysis were: 
[CuI(NHC2)]+ = MO52a-HOMO; [CuII(NHC2)]2+ = MO60a-(SOMO); [CuIII(NHC2)]3+ = MO60a-(HOMO); 
[CuIII(NHC4)]3+ = 59a-HOMO; [Cu3+(CF3)4]− = 42a-HOMO. 
 
 
Figure S25 (Cu) 3dxz–2px (C) b2 σ-type MO in [CuI(NHC2)]+ as viewed along the principle C2 axis (z). Isosurfaces 
are plotted at a value of 0.05 a.u. 


--- Page 26 ---
 
26 
 
Figure S26 Calculated total and x,y,z polarized VtC spectra for the complexes investigated in the present work. 
The x,y and z axes are as depicted in Scheme 1. 
Due to the relatively high symmetry and approximate planar structure of the [Cun+(NHC2/4)]n+ complexes, the 
calculated polarized VtC XES plots are highly informative to understand various intensity contributions. Figure 
S26 shows that the majority, if not all the intensity for the [Cun+(NHC2)]n+ series is x-polarized, corresponding to 
the Cu-CNHC vector. Each of the [Cun+(NHC2)]n+ complexes exhibit a dominant peak along the x-direction of 
comparable intensity, although [CuII(NHC2)]2+ is noticeably the least intense as well as being broader due to the 
distribution of transitions (Figures 5 and S24). Despite the decreasing average CNHC-Cu bond lengths from the 
Cu(I) species (1.94 Å) to the Cu(II) (1.91 Å) and Cu(III) (1.88 Å), there is not a clear trend of integrated intensity. 
However, inspection of the C-Cu-C angles in these complexes shows that the angle in the Cu(I) species (173.9°) 
and the Cu(III) species (172.7°) are closer to linear than that in the Cu(II) species (169.0°). The more linear C-
Cu-C σ bonds in the Cu(I) and Cu(III) species will have better overlap with the Cu npx orbitals, thus resulting in 
larger intensity being imparted into the transition. This shows that the decreasing Cu-C bond lengths in the Cu(II) 
species relative to the Cu(I) species are counteracted by the geometry changes resulting from the oxidation of 
the Cu center. Such an angular dependence on the VtC spectrum of hydroxo and oxo-bridged Fe dimers, as well 
as Ni NOx complexes has been demonstrated previously.20,21 
In the orthogonal y-direction, along the Cu-Npyridine vector, the calculated VtC XES exhibits minimal intensity for 
the [CuI(NHC2)]+ complex, consistent with the pyridines being assigned as non-coordinating. However, for 
[CuII(NHC2)]2+ and [CuIII(NHC2)]3+, an increasing amount of intensity is observed in the y-polarized direction due 
to the coordinating pyridines (Figure S26). Inspection of the MOs corresponding to the most dominant transitions 
qualitatively reveals localized pyridine character. Furthermore, inspection of the Löwdin population analysis of 
the MOs corresponding to the transitions in this shoulder region reveals increased pyridine p-orbital character 
relative to the Cu(I) species (Table S9). Lastly, both [CuII(NHC2)]2+ and [CuIII(NHC2)]3+ exhibit a small amount of 
z-polarized intensity that may also be attributed to either the weakly coordinated acetonitrile ligand or a slight 
geometric distortion (Figures S22 and S26). 


--- Page 27 ---
 
27 
 
 
Figure S27 Summation of the emission intensity in three main regions of interest in the calculated VtC spectra 
for [CuI(NHC2)]+ (top), [CuII(NHC2)]2+ (middle) and [CuIII(NHC2)]3+ (bottom). MO isosurfaces are plotted at a value 
of 0.05 a.u. 
 
 
 
 
 
 


--- Page 28 ---
 
28 
Table S9 Löwdin populations of canonical orbitals contributing to the calculated emission intensity of the 
[Cun+(NHC2)]n+ series in three different regions of interest as shown in Figure S27. 
 
Cu d,p [%] Total ligand s,p [%] NHC s,p [%] 
Pyridine s,p [%] Acetonitrile s,p [%] 
[CuI(NHC2)]+ 
 
 
 
 
 
Region 1 
0, 0 
53.0, 39.1 
34.0, 27.5 
14.2, 8.9 
- 
Region 2 
0, 0.6 
28.1, 66.6 
14.4, 40.1 
5.1, 16.3 
- 
Region 3 
9.3, 7.9 
9.8, 64.6 
8.0, 40.9 
0.8, 21.9 
- 
a[CuII(NHC2)]2+ 
 
 
 
 
 
Region 1 
0.1, 1.0 
22.9, 73.3 
11.9, 33.9 
5.2, 26.7 
1.0, 7.6 
Region 2 
17.7, 0.8 
13.7, 60.6 
1.9, 15.6 
4.6, 20.4 
2.2, 12.5 
Region 3 
16.7, 4.1 
7.1, 63.1 
2.8, 21.9 
1.7, 27.5 
1.4, 9.3 
[CuIII(NHC2)]3+ 
 
 
 
 
 
Region 1 
0.4, 1.4 
20.4, 75.1 
9.2, 31.5 
6.2, 29.7 
0.2, 0.4 
Region 2 
18.3, 0.6 
17.4, 57.6 
3.3, 18.6 
4.9, 19.7 
2.1, 7.2 
Region 3 
4.9, 4.5 
9.4, 72.0 
2.8, 18.7 
1.9, 20.7 
4.2, 30.8 
a Values for the [CuII(NHC2)]2+ species are the average for the α and β spin manifolds. 
 
 


--- Page 29 ---
 
29 
 
Figure S28 DFT calculated quadrupole allowed only VtC emission spectra and respective IWAEs for all 
complexes in the present analysis. A Gaussian line broadening function was used and a FWHM value of 0.8 eV 
was selected.  
 
Table S10 B3LYP/def2-TZVP calculated total Cu nd orbital admixture into valence MOs (and total Cu spin 
population) given by Löwdin and Mulliken population analysis of either the canonical orbitals, unrestricted natural 
orbitals (UNOs) or quasi restricted orbitals (QROs).  
 
% Cu nd admixture 
Cu spin population [%] 
 
Löwdin  
canonical (UNO) [QRO] 
Mulliken  
canonical (UNO) [QRO] 
Löwdin (Mulliken) 
canonical 
[CuI(NHC2)]+  
HOMO 
58.3 (1.1) [58.3] 
57.4 (0.4) [57.5] 
- 
[CuII(NHC2)]2+  
HOMO 
29.5 (61.0) [61.0] 
27.3 (59.8) [60.0] 
61.5 (61.7) 
[CuIII(NHC2)]3+  
LUMO 
42.0 (2.4) [42.0] 
37.3 (1.6) [37.3] 
- 
[CuIII(NHC4)]3+ 
LUMO 
45.2 (0.6) [47.4] 
33.8 (0.6) [35.2] 
- 
[CuIII(CF3)4]– 
LUMO 
38.8 (7.3) [38.0] 
24.1 (7.6) [24.4] 
- 
 
D4h [CuCl4]2–  
HOMO 
20.2 (58.9) [59.0] 
15.5 (56.6) [56.6] 
55.6 (51.5) 
Values are given for the α orbitals.  
 


--- Page 30 ---
 
30 
UV-Vis  
 
Figure S29 TDDFT calculated (B3LYP/ZORA-def2-TZVP/ZORA) UV-Vis spectrum for [Cu(CF3)4]− using the 
conductor-like polarizable continuum model (ε = inf.) to account for the negative charge of the complex anion. 
Donor and acceptor MOs are plotted for selected transitions along with their weighted contributions. Isosurfaces 
are plotted at a value of 0.05 a.u. 
 


--- Page 31 ---
 
31 
 
Figure S30 TDDFT calculated (B3LYP/ZORA-def2-TZVP/ZORA) UV-Vis spectrum for [Cu(CF3)4]− and 
[CuIII(NHC4)]3+ using the conductor-like polarizable continuum model (ε = inf.) to account for the overall charge of 
the complex ions. Donor and acceptor MOs are plotted for selected transitions along with their weighted 
contributions. Isosurfaces are plotted at a value of 0.05 a.u. 
 
 
 


--- Page 32 ---
 
32 
 
Example Orca Input Files 
Geometry Optimization 
! UKS RIJCOSX B3LYP def2-TZVP AutoAux opt TightSCF Grid4 FinalGrid7 SlowConv NormalPrint  
! UNO CPCM(if complex carries non-zero charge)  
 
%pal 
nprocs 12 
end 
 
%maxcore 3000 
 
%scf 
MaxIter 700 
End 
 
*xyzfile charge spin-multiplicity xyz_filename.xyz 
 
 
X-ray Absorption Spectroscopy 
! UKS RIJCOSX B3LYP ZORA ZORA-def2-TZVP AutoAux TightSCF SlowConv Grid4 FinalGrid7 Pal8 LargePrint 
! CPCM UNO 
  
%scf 
Maxiter 700 
end 
 
%maxcore 3000 
  
%tddft   
NRoots 200   
MaxDim 900     
OrbWin[0] = 0, 0, -1, -1   
OrbWin[1] = 0, 0, -1, -1  
Doquad true   
end  
  
*xyzfile charge spin_multiplicity xyz_filename.xyz 
X-ray Emission Spectroscopy 
! UKS RIJCOSX B3LYP ZORA ZORA-def2-TZVP AutoAux Grid4 FinalGrid7 SlowConv TightSCF pal8  
! LargePrint UNO CPCM(if complex carries non-zero charge) 
 
%scf 
MaxIter 700 
end 
 
%maxcore 3000 
 
%xes 


--- Page 33 ---
 
33 
CoreOrb 0,0 
OrbOp 0,1 
CoreOrbSOC 0,1 
DoSOC true 
end 
 
*xyzfile charge spin_multiplicity xyz_filename.xyz 
 
Löwdin and Mulliken Population analysis of QROs 
! LargePrint NoIter MOREAD def2-TZVP 
 
%moinp "tddft_calculation_filename.qro" 
 
*xyzfile charge spin_multiplicity xyz_filename.xyz 
 
Broken Symmetry 
! UKS RIJCOSX B3LYP def2-TZVP AutoAux Grid4 FinalGrid7 SlowConv TightSCF pal8  
! LargePrint UNO CPCM 
 
%scf 
BrokenSym 1,1 
end 
 
%maxcore 3000 
 
*xyzfile charge triplet_spin_multiplicity xyz_filename.xyz 
 
UV-vis spectroscopy 
! UKS RIJCOSX B3LYP ZORA ZORA-def2-TZVP AutoAux Grid4 FinalGrid7 SlowConv  
! TightSCF pal8 NormalPrint UNO CPCM 
 
%scf 
MaxIter 700 
end 
 
%maxcore 3000 
 
%tddft 
NRoots 60 
MaxDim 5 
end 
 
*xyzfile charge triplet_spin_multiplicity xyz_filename.xyz 
 
 


--- Page 34 ---
 
34 
XYZ Coordinates for Orca Calculations 
[CuI(NHC2)]+ 
Charge = 1+, Spin Multiplicity = 1 
  Cu  -0.00740413299260     -0.04079336237917      0.13207803040467 
  N   -0.02306510512784     -2.62451096464465      0.39810987921352 
  N   0.00932859614757      2.56090157725132      0.28689515106662 
  C   1.94728414746307     -0.07464419591113      0.22304221810076 
  C   -1.96070678328296      0.00365151709803      0.24885425026110 
  C   1.08157935414333     -2.80336943362877      1.12738717334633 
  C   -1.21411673613655     -2.79297614879768      0.97584022771510 
  C   -1.08332126019588      2.76462195233140      1.02745388054978 
  C   1.20918221142673      2.74988236796822      0.83932751055978 
  N   2.77476254700957     -1.13712084620435      0.39209544958973 
  N   2.79304105107554      0.98787826466697      0.18784269586343 
  N   -2.78672815754387      1.07280132298858      0.37910296228816 
  N   -2.80684976197257     -1.05915199921355      0.26438925445634 
  C   2.37810980747925     -2.55535313027989      0.40701562104720 
  C   1.03649430977973     -3.16881219046406      2.46758777732075 
  C   -1.35064324389426     -3.17142403630490      2.30676528152688 
  C   -2.40085380900937     -2.46631470888475      0.11221353571653 
  C   -2.39162657503318      2.49119399335755      0.33812444381695 
  C   -1.01655894519240      3.17624761781704      2.35313538783983 
  C   1.36735837672303      3.17618543289052      2.15318018745124 
  C   2.38373841133867      2.38632893324904     -0.02548123351970 
  C   4.10057715692895     -0.74861859514705      0.44610284290877 
  C   4.11180015935928      0.59674217931508      0.31921745065780 
  C   -4.11228945727196      0.68712224807323      0.45745623321823 
  C   -4.12467488342361     -0.66230320389551      0.38682817435480 
  C   -0.20365214973805     -3.35912182510847      3.06403240449292 
  C   0.23298074223064      3.38932364419960      2.92192097796087 
  H   3.18870313031190     -3.10199987409391      0.88468436480272 
  H   2.28322449815311     -2.90045117229094     -0.62234805289014 
  H   1.94963117983410     -3.31241259409035      3.02998583768353 
  H   -2.33170808456379     -3.31416484129875      2.74082452252246 
  H   -2.14975535790919     -2.65555525230473     -0.93040014891398 
  H   -3.26187618230452     -3.07406076961440      0.38354869009515 
  H   -3.19455076072294      3.05305206474972      0.81098475814832 
  H   -2.31558271709328      2.80123440237509     -0.70399280345024 
  H   -1.92058719760087      3.33742714306229      2.92528875481498 
  H   2.35541366793588      3.33417698636880      2.56562933115066 
  H   2.11753966469871      2.52565130381315     -1.07220997868685 
  H   3.24760556566286      3.00677083226666      0.20491660408108 


--- Page 35 ---
 
35 
  H   4.90492510117534     -1.45075645416561      0.57301593914522 
  H   4.92890452820429      1.29540084567437      0.31498690058940 
  H   -4.91574059417615      1.39410577352380      0.56084134748227 
  H   -4.94171318776417     -1.36039101945015      0.41814394547521 
  H   -0.27366218218347     -3.65684137259083      4.10188738919502 
  H   0.31966790105189      3.72258782672313      3.94755548354672 
 
[CuII(NHC2)]2+ 
Charge = 2+, Spin Multiplicity = 2 
  Cu  0.10082553881777     -0.01162902283546      0.01361636182430 
  N   0.09177864474464     -2.19517970873212      0.32308094264054 
  N   0.09137194740309      2.19001728077502      0.10570309438298 
  N   0.10827421165350     -0.07099431073676     -2.35899804452316 
  C   2.01124054044758     -0.01428398467647      0.28199763732638 
  C   -1.82323092576793      0.01339166647361      0.11270973440238 
  C   1.18034605791675     -2.91675405129246     -0.00751777317502 
  C   -0.99026703699657     -2.83143876581840      0.80973520770289 
  C   -0.93174111333120      2.86976818341817     -0.44209041501289 
  C   1.14371215887214      2.87334360283847      0.59769502669000 
  C   -0.11718869170953      0.02223724123049     -3.48195327172126 
  N   2.84037624779476     -1.02856194168101      0.00799105287700 
  N   2.77700097053876      0.95431556545254      0.80012869398815 
  N   -2.61612552028967      0.99388483820824     -0.34147234682676 
  N   -2.64373123466321     -0.91174734175171      0.62498812702716 
  C   2.38269983086823     -2.24831526300419     -0.64617296499871 
  C   1.21761556807379     -4.29737019646537      0.14632964023301 
  C   -1.02253367749143     -4.21586666791314      0.94268663194138 
  C   -2.17909959104891     -2.07492271707419      1.36978218695061 
  C   -2.10415391943216      2.14384708407699     -1.07304378293417 
  C   -0.92461996202856      4.25766057168222     -0.52933508712561 
  C   1.21860276702133      4.25803486249495      0.50658027466162 
  C   2.22929371133160      2.17144489220797      1.38788909987281 
  C   -0.41604792704665      0.15140308518597     -4.89588906257067 
  C   4.14517063482166     -0.70555602282581      0.35413558928361 
  C   4.10388707112358      0.54935114347680      0.86280404025582 
  C   -3.95189824633246      0.68780182495195     -0.12024371896655 
  C   -3.96936544630035     -0.51522370940242      0.50154746212983 
  C   0.09689507503996     -4.96063801549148      0.61736188638752 
  C   0.17302909058302      4.96230233947447     -0.06662095740645 
  H   3.21066160351901     -2.95194856739489     -0.64678144087105 
  H   2.13895157561739     -2.00919453198113     -1.68446866059659 
  H   2.11412694080692     -4.84477824989436     -0.11005708341445 
  H   -1.91335528157941     -4.69984695798052      1.31905774454038 


--- Page 36 ---
 
36 
  H   -3.01551588041271     -2.76487493900858      1.44658839809787 
  H   -1.92902229617680     -1.75297190112432      2.38566071904375 
  H   -2.92268887729186      2.84834785074471     -1.19330006104934 
  H   -1.80378410783041      1.80682969024233     -2.06775743916666 
  H   -1.76878088519060      4.77563038127525     -0.96352176834137 
  H   2.08386636337821      4.77757918354921      0.89433559223726 
  H   3.05339582206838      2.86444222148538      1.53481952833578 
  H   1.83103037585415      1.92599991430796      2.37701918926787 
  H   -0.78462089833580     -0.79806730425287     -5.28657247832320 
  H   0.48137764796439      0.44315832432087     -5.44263359472435 
  H   -1.18177964697345      0.91317245703750     -5.04664450917573 
  H   4.97275893138423     -1.37640146657256      0.20585017751783 
  H   4.88667044492290      1.17136807304030      1.25885312452422 
  H   -4.75715137086282      1.33490899090250     -0.41987965305984 
  H   -4.79145713149439     -1.10545798812091      0.86591334386195 
  H   0.09857593095695     -6.03605126121241      0.73575290932416 
  H   0.21004011506214      6.04089753838973     -0.14160812734517 
 
[CuIII(NHC2)]3+ 
Charge = 3+, Spin Multiplicity = 1 
  Cu  -0.07557324534937      0.08473907599887      0.01749033748120 
  N   -0.11710427142914     -1.92599544467153      0.18499604266301 
  N   -0.03408855310710      2.09834319628169      0.15192748564677 
  N   -0.06374217098734      0.01037443838644     -2.35450306979143 
  C   -1.96331455466345      0.09258358873240      0.11688310543760 
  C   1.81096597131832      0.07711566686513      0.16051460223402 
  N   -2.77232298670605     -0.82451915396435     -0.40362265527385 
  N   -2.72203492841614      0.98952662489733      0.73832089974582 
  N   2.63204382830595      0.97882855247198     -0.36652835933578 
  N   2.55228490983997     -0.79376132700535      0.83765055387342 
  C   0.87877312657885     -2.59424048389109      0.82659624639028 
  C   -1.09389580910824     -2.63335378563464     -0.43741378792053 
  C   -4.09293111422718     -0.50350057690443     -0.11784575943584 
  C   -2.25768058060771     -1.96631986720113     -1.13886379999260 
  C   -4.06096674973948      0.64521354291859      0.60729739372831 
  C   -2.13349050179208      2.06810777191251      1.52145382510592 
  C   -1.02672185965424      2.78103419215742      0.78058873325227 
  C   0.97591682431046      2.79223106044453     -0.43609813244297 
  C   2.11844492599363      2.09608390253589     -1.14279545936202 
  C   3.94375593143155      0.67702411647537     -0.02512379657418 
  C   3.89331171515817     -0.44723814095791      0.73785917581180 
  C   1.93263293339283     -1.85756569651206      1.61918072477104 
  C   -0.09402854849986     -0.10324945645329     -3.49935864804762 


--- Page 37 ---
 
37 
  C   0.93675857382873     -3.97870713513175      0.81292299897897 
  C   -1.05987802315533     -4.02119018052970     -0.47321905287160 
  C   -1.03533297770752      4.16774288692284      0.81203754717493 
  C   0.99323556893107      4.17966249677476     -0.42356411948387 
  C   -0.14275362167436     -0.25197924607384     -4.93928720316056 
  C   -0.02976206203020     -4.70955078177724      0.14198565773654 
  C   -0.02058348487005      4.88439121350562      0.20157150153427 
  H   -4.92310351946544     -1.10006077819544     -0.45629818867217 
  H   -1.95560604924565     -1.63971717101380     -2.13446042588483 
  H   -3.06023219546215     -2.69003245669721     -1.24701954315383 
  H   -4.85678960071782      1.22669557552605      1.04110113890439 
  H   -1.75359426036551      1.65753061436998      2.46168888551861 
  H   -2.91483345337440      2.78007962827257      1.76997881816347 
  H   1.79037051401078      1.72945201992668     -2.11654112956889 
  H   2.92484818011919      2.80617806828481     -1.29933571930475 
  H   4.78209301368428      1.26733943187705     -0.35442103291303 
  H   4.67697464507442     -1.00891539796587      1.21761859523896 
  H   1.49438532830543     -1.42573816220245      2.52335930794152 
  H   2.70413584773577     -2.55494548654925      1.93135083606764 
  H   1.74464354535760     -4.47883796986312      1.32916578102452 
  H   -1.85262374946789     -4.55439654421655     -0.98058415894477 
  H   -1.84120594589389      4.67892136199888      1.32084502853338 
  H   1.81370984554407      4.70073799760100     -0.89806255753123 
  H   -0.88708120023241     -1.00215103290511     -5.21138561163383 
  H   0.83174306785771     -0.57077220076510     -5.31418354922233 
  H   -0.41730287074995      0.69834839761058     -5.40116043977222 
  H   0.01265717110420     -5.79086753088648      0.11307512896297 
  H   -0.01545065818304      5.96677143621972      0.22140529737308 
 
 


--- Page 38 ---
 
38 
[CuIII(NHC4)]3+ 
Charge = 3+, Spin Multiplicity = 1 
  Cu  0.03229808776766     -0.00802812810694      0.11864237027951 
  N   -2.72236441013310     -1.01822590913991     -0.14960253623219 
  N   -2.70775966645164      1.05165990371497      0.36015140049886 
  N   -0.99182677179212      2.73663661758113      0.42799817821532 
  N   1.10074448186391      2.73600010808165      0.01540624132707 
  N   2.78810731813870      1.01965621725733     -0.04001317343073 
  N   2.76758880744168     -1.08680411723597      0.29827934168672 
  N   1.05070082187890     -2.77168505299039      0.22578239216403 
  N   -1.03708489976242     -2.73351942313217     -0.20703421110259 
  C   -1.89907920868679      0.01048829559839      0.10758493656803 
  C   -4.05775312012779     -0.63485251059755     -0.05570934873197 
  C   -4.04869988195395      0.68771926887409      0.26135656080057 
  C   -5.19138504167512     -1.57367244889198     -0.25897594410857 
  H   -6.12600910882298     -1.06898772618353     -0.02356893674340 
  H   -5.10536620157276     -2.43969702346000      0.40068144816933 
  H   -5.24528614831799     -1.93546636579990     -1.28849847318689 
  C   -5.17009958081677      1.64258438237474      0.45740994034379 
  H   -6.11029521854440      1.15027837403222      0.21767312335027 
  H   -5.06873379092374      2.50662916473243     -0.20273830230252 
  H   -5.22470287823333      2.00610692132287      1.48623488531919 
  C   -2.31219154250611      2.34561998129228      0.87111040795094 
  H   -2.33709253692771      2.32833531268671      1.96202245862370 
  H   -3.02204397616674      3.08265205266051      0.51020289364037 
  C   0.04929811497074      1.92107752399187      0.19712738025223 
  C   -0.60614421139231      4.07400698785271      0.38639837453204 
  C   0.72966576443080      4.07370170489001      0.13124650642465 
  C   -1.55511458342522      5.20210040963273      0.57063851288606 
  H   -1.04222487275064      6.14195952138863      0.37751493107682 
  H   -1.96282940532215      5.23351241760329      1.58376877477134 
  H   -2.39064239485964      5.12937486658476     -0.12870112140424 
  C   1.68813126636855      5.20216727639866      0.00619637728472 
  H   1.18240771468055      6.13518675668267      0.24578176616021 
  H   2.09832373591545      5.28245937720486     -1.00309526977729 
  H   2.52151920070277      5.08699086027251      0.70269328150244 
  C   2.41612631650466      2.35588877618411     -0.45226475926732 
  H   2.44440679944015      2.41932726919842     -1.54145435455628 
  H   3.13762346138938      3.05191888289775     -0.03670258513487 
  C   1.96242168578846     -0.02585449938165      0.12790937212463 
  C   4.12236796040711      0.62549381989688      0.02620426476268 
  C   4.10933926116591     -0.71846057305062      0.23343204507798 
  C   5.26123619714263      1.57192141960304     -0.09399483575030 


--- Page 39 ---
 
39 
  H   6.19205370099001      1.04611167414192      0.10763270649858 
  H   5.17415179500235      2.38419019165571      0.63054019690890 
  H   5.32398484517666      2.01323409360016     -1.09136045474970 
  C   5.22958594875976     -1.68700759949008      0.35361922666727 
  H   6.17097695510800     -1.17776251116382      0.15862651447087 
  H   5.13053328307037     -2.49341065413460     -0.37621741315490 
  H   5.27957754154770     -2.13449804127892      1.34893919740984 
  C   2.36828377879116     -2.41610887998863      0.70653077817220 
  H   2.38602835536991     -2.48132685080366      1.79578885642928 
  H   3.08034367616995     -3.12447702349841      0.29596120001931 
  C   0.01598945350023     -1.93740482798923      0.03696398141756 
  C   0.65782363088973     -4.10203986401947      0.09946190899722 
  C   -0.67541172759528     -4.07776919045526     -0.16798302768235 
  C   1.59657171863375     -5.24699527430679      0.22480226222433 
  H   1.07710729473792     -6.17042463433118     -0.02226232696780 
  H   1.99983599748448     -5.33883901388458      1.23589825924429 
  H   2.43561699786388     -5.14185993409134     -0.46630372824764 
  C   -1.64163317753147     -5.18853478484637     -0.36743030588269 
  H   -1.14780280834696     -6.13731469890547     -0.16818145322813 
  H   -2.03474670552326     -5.21195553448894     -1.38662598982001 
  H   -2.48586625839380     -5.10223468886956      0.31952199784533 
  C   -2.34451243479010     -2.31620510404550     -0.66324650498473 
  H   -3.07156508050730     -3.04201412473838     -0.31438641450005 
  H   -2.35692955924031     -2.29174160558918     -1.75433040715131 
 
[CuIII(CF3)4]− 
Charge = 1-, Spin Multiplicity = 1 
  C   -0.57888321333209     -0.87192810770068     -4.61669655249915 
  C   1.24870669476421     -3.05020976763769     -4.40963955185390 
  C   1.47833807231003      0.87018718488180     -3.72877744788674 
  C   3.37418382159552     -1.11507810433457     -4.53219278742627 
  Cu  1.38496409074348     -1.05037661636376     -4.31751573386117 
  F   3.87088529843570      0.02400621230627     -5.10562741861118 
  F   3.79458304414789     -2.10584524235780     -5.37497292534176 
  F   4.08208039740395     -1.28693565968239     -3.38227260075224 
  F   2.51557920676023      1.11013329617410     -2.87255884171744 
  F   1.61325680072138      1.77897519107215     -4.73225108258272 
  F   0.37859040641269      1.26585090019746     -3.01888290997497 
  F   -1.08831195537215     -1.80859502505455     -5.47048544593969 
  F   -0.89399280688326      0.31876306969050     -5.21409482456034 
  F   -1.34336207842653     -0.94156858200232     -3.49436571889668 
  F   0.13488502680767     -3.54372229366747     -3.78682883279539 
  F   1.22150934303065     -3.56228088779489     -5.66753459725558 


--- Page 40 ---
 
40 
  F   2.27957453288055     -3.68768122272611     -3.77795042304449 
 
References  
(1)  
Liu, Y.; Resch, S. G.; Klawitter, I.; Cutsail III, G. E.; Demeshko, S.; Dechert, S.; Kühn, F. E.; DeBeer, S.; 
Meyer, F. An Adaptable N-Heterocyclic Carbene Macrocycle Hosting Copper in Three Oxidation States. 
Angew. Chemie Int. Ed. 2020, 59 (14), 5696–5705. https://doi.org/10.1002/anie.201912745. 
(2)  
Romine, A. M.; Nebra, N.; Konovalov, A. I.; Martin, E.; Benet-Buchholz, J.; Grushin, V. V. Easy Access to 
the Copper(III) Anion [Cu(CF )4]−. Angew. Chemie Int. Ed. 2015, 54 (9), 2745–2749. 
https://doi.org/10.1002/anie.201411348. 
(3)  
Bernd, M. A.; Dyckhoff, F.; Hofmann, B. J.; Böth, A. D.; Schlagintweit, J. F.; Oberkofler, J.; Reich, R. M.; 
Kühn, F. E. Tuning the Electronic Properties of Tetradentate Iron-NHC Complexes: Towards Stable and 
Selective Epoxidation Catalysts. J. Catal. 2020, 391, 548–561. https://doi.org/10.1016/j.jcat.2020.08.037. 
(4)  
Ghavami, Z. S.; Anneser, M. R.; Kaiser, F.; Altmann, P. J.; Hofmann, B. J.; Schlagintweit, J. F.; Grivani, 
G.; Kühn, F. E. A Bench Stable Formal Cu(III): N-Heterocyclic Carbene Accessible from Simple Copper(II) 
Acetate. Chem. Sci. 2018, 9 (43), 8307–8314. https://doi.org/10.1039/c8sc01834k. 
(5)  
Sheldrick, G. M. SHELXT - Integrated Space-Group and Crystal-Structure Determination. Acta 
Crystallogr. Sect. A Found. Crystallogr. 2015, 71 (1), 3–8. https://doi.org/10.1107/S2053273314026370. 
(6)  
SADABS; Bruker AXS GmbH. Karlsruhe, Germany 2016. 
(7)  
Ravel, B.; Newville, M. ATHENA, ARTEMIS, HEPHAESTUS: Data Analysis for X-Ray Absorption 
Spectroscopy 
Using 
IFEFFIT. 
J. 
Synchrotron 
Radiat. 
2005, 
12 
(4), 
537–541. 
https://doi.org/10.1107/S0909049505012719. 
(8)  
Neese, F. The ORCA Program System. Wiley Interdiscip. Rev. Comput. Mol. Sci. 2012, 2 (1), 73–78. 
https://doi.org/10.1002/wcms.81. 
(9)  
Neese, F. Software Update: The ORCA Program System, Version 4.0. Wiley Interdiscip. Rev. Comput. 
Mol. Sci. 2018, 8 (1), 4–9. https://doi.org/10.1002/wcms.1327. 
(10)  
Lee, C.; Yang, W.; Parr, R. G. Development of the Colic-Salvetti Correlation-Energy Formula into a 
Functional of the Electron Density. Phys. Rev. B 1988, 37 (2), 785–788. 
(11)  
Becke, A. D. Density-Functional Thermochemistry. III. The Role of Exact Exchange. J. Chem. Phys. 1993, 
98 (7), 5648–5652. https://doi.org/10.1063/1.464913. 
(12)  
Weigend, F.; Ahlrichs, R. Balanced Basis Sets of Split Valence, Triple Zeta Valence and Quadruple Zeta 
Valence Quality for H to Rn: Design and Assessment of Accuracy. Phys. Chem. Chem. Phys. 2005, 7 
(18), 3297–3305. https://doi.org/10.1039/b508541a. 
(13)  
Stoychev, G. L.; Auer, A. A.; Neese, F. Automatic Generation of Auxiliary Basis Sets. J. Chem. Theory 
Comput. 2017, 13 (2), 554–562. https://doi.org/10.1021/acs.jctc.6b01041. 
(14)  
Neese, F.; Wennmohs, F.; Hansen, A.; Becker, U. Efficient, Approximate and Parallel Hartree-Fock and 
Hybrid DFT Calculations. A “chain-of-Spheres” Algorithm for the Hartree-Fock Exchange. Chem. Phys. 


--- Page 41 ---
 
41 
2009, 356 (1–3), 98–109. https://doi.org/10.1016/j.chemphys.2008.10.036. 
(15)  
Klamt, A.; Schüürmann, G. COSMO: A New Approach to Dielectric Screening in Solvents with Explicit 
Expressions for the Screening Energy and Its Gradient. J. Chem. Soc. Perkin Trans. 2 1993, No. 5, 799–
805. https://doi.org/10.1039/P29930000799. 
(16)  
Van Lenthe, E.; Wormer, P. E. S.; Van Der Avoird, A. Density Functional Calculations of Molecular G-
Tensors in the Zero-Order Regular Approximation for Relativistic Effects. J. Chem. Phys. 1997, 107 (7), 
2488–2498. https://doi.org/10.1063/1.474590. 
(17)  
Pantazis, D. A.; Chen, X. Y.; Landis, C. R.; Neese, F. All-Electron Scalar Relativistic Basis Sets for Third-
Row 
Transition 
Metal 
Atoms. 
J. 
Chem. 
Theory 
Comput. 
2008, 
4 
(6), 
908–919. 
https://doi.org/10.1021/ct800047t. 
(18)  
Gasparro, F. P.; Kolodny, N. H. NMR Determination of the Rotational Barrier in N,N-Dimethylacetamide. 
A Physical Chemistry Experiment. J. Chem. Educ. 1977, 54 (4), 258. https://doi.org/10.1021/ed054p258. 
(19)  
Bill, E. JulX, Program for Simulation of Molecular Magnetic Data. Max-Planck Institute for Chemical 
Energy Conversion: Mülheim/Ruhr 2008. 
(20)  
Pollock, C. J.; Lancaster, K. M.; Finkelstein, K. D.; DeBeer, S. Study of Iron Dimers Reveals Angular 
Dependence of Valence-to-Core X-Ray Emission Spectra. Inorg. Chem. 2014, 53 (19), 10378–10385. 
https://doi.org/10.1021/ic501462y. 
(21)  
Phu, P. N.; Gutierrez, C. E.; Kundu, S.; Sokaras, D.; Kroll, T.; Warren, T. H.; Stieber, S. C. E. Quantification 
of Ni-N-O Bond Angles and NO Activation by X-Ray Emission Spectroscopy. Inorg. Chem. 2021, 60 (2), 
736–744. https://doi.org/10.1021/acs.inorgchem.0c02724.